\documentclass[aspectratio=169, t, 10pt]{beamer}
\usetheme{Boadilla}
\beamertemplatenavigationsymbolsempty

\usepackage{standalone}
\usepackage[utf8]{inputenc}
\usepackage[english]{babel}
\usepackage{pgfplots}
\usepackage{mathtools}
\mathtoolsset{showonlyrefs}
\usepackage{qrcode}
\usepackage{multimedia}

\newcommand\Mean[1]{\mathbb{E}\!\left[#1\right]}
\newcommand\Var[1]{\mathbb{V}\!\left[#1\right]}
\newcommand\Cov[2]{\mathrm{Cov}\!\left[#1,#2\right]}
\newcommand\Gauss[2]{\mathcal{N}\!\left({#1},\,{#2}\right)}
\newcommand\GP[2]{\mathcal{GP}\!\left({#1},\,{#2}\right)}
\newcommand{\Identity}{\mathbb{I}}

\title[Correlation based travel time inversion]{Bayesian Travel Time Inversion adopting Gaussian Process Regression}
\subtitle{-- with the focus on uncertainty analysis --}
\author[\tt mauerber@uni-potsdam.de]{Stefan Mauerberger \and Matthias Holschneider}
\institute[Math@UP]{University Potsdam, Institute of Mathematics}
\titlegraphic{%\vspace{-1cm}
              \hspace{2cm}
              \parbox[c]{0.17\linewidth}{\includegraphics[width=\linewidth]{./logos/GeoSim_Logo}}
              \hfill
              \parbox[c]{0.10\linewidth}{\includegraphics[width=\linewidth]{./logos/UniPotsdam_Logo}}
              \hspace{2cm}
              }
\date[AGU~2017]{AGU Fall Meeting -- 13\textsuperscript{th} December 2017}

\begin{document}
\input{def_example}

\frame[noframenumbering, plain]
    {\maketitle}

\begin{frame}
    \frametitle{Overview}
    \framesubtitle{Where are we driving at -- What I am going to present}

\begin{columns}
\column{.55\textwidth}

    \begin{block}{Surface Wave Tomography}
        \begin{itemize}
            \item Correlation-based Bayesian inversion
            \item with ambient noise records
            \item Explicit correlation structure
            \item Observational functional from observational functional
        \end{itemize}
        \hfill {\Large $\leadsto$} Realistic, location dependent model uncertainties ~
    \end{block}

    \begin{exampleblock}{First steps}
        \begin{itemize}
            \item A synthetic test for the static field
            \item Explicit spatial correlation structure
            \item Proxy correlations for declination, inclination and intensity
        \end{itemize}
        \hfill {\Large $\leadsto$} A\,posteriori distribution for the static field ~
    \end{exampleblock}


    Simplifications
    \begin{itemize}
        \item Synthetic test
        \item Proof of concept
        \item No dispersion
        \item No refraction
    \end{itemize}

    Setting
    \begin{itemize}
        \item Subset of ScanArray \SFWnst
        \item Two large scale anomalies
        \item Pseudo observations
    \end{itemize}


\column[T]{.44\textwidth}
    \vspace{-10mm}
    \input{fig_reference_model.pgf}
\end{columns}

\end{frame}

\begin{frame}
    \frametitle{Travel Time Observations}
    \framesubtitle{a functional w.r.t.~an underlying velocity model}

\begin{columns}
\column{.55\textwidth}%
    \begin{equation}
        d = \mathrm T_{s,r}[v] + \varepsilon = \int_C \frac 1{v(r)} \mathrm d r + \varepsilon
    \end{equation}
    \begin{description}[leftmargin=! ,labelwidth=1cm]
        \item [Measured value]           $d$
        \item [Source location]          $s$
        \item [Receiver position]        $r$
        \item [Observational functional] $\mathrm T_{s,r}[\cdot]$
        \item [Velocity model]           $v$
        \item [Ray path]                 $C$
        \item [Measurement error]        $\varepsilon$
    \end{description}

    \begin{alertblock}{It is the velocity model what we are after}
    Travel times are putting an Integral-constraint on $v$
    \end{alertblock}

\column[T]{.44\textwidth}
    \vspace{-10mm}
    \input{fig_path_coverage.pgf}

\end{columns}

\end{frame}


\begin{frame}
    \frametitle{A\,Priori Velocity Model}
    \framesubtitle{expressing our ignorance}

\begin{columns}
\column{.55\textwidth}%
    \begin{block}{Model $v(x)$ a Gaussian Random Field}
    \begin{equation}
        v \to V \sim \GP{\mu_V}{K_V}
    \end{equation}
    \begin{description}[leftmargin=! ,labelwidth=6cm]
        \item [Prior mean function] $\mu_V(x) = \SFWmuCpri \, \frac ms$
        \item [Covariance function] $K_V(x,y)$
    \end{description}
    \end{block}
    \medskip

    Correlation Kernel
    \begin{equation}
        K_V(x,y) = \tau^2 \exp\left\{ -\frac12 \frac{d(x,y)^2}{\ell^2}\right\}
    \end{equation}
    \begin{description}[leftmargin=! ,labelwidth=6cm]
        \item [Standard deviation]   $\tau = \SFWtau \, \frac ms$
        \item [Great circle distance] $d(x,y)$
        \item [Characteristic lenth]  $\ell = \SFWell \, m$
    \end{description}

\column[T]{.44\textwidth}
    \vspace{-10mm}
    \input{fig_kernel.pgf}

\end{columns}

\end{frame}


\begin{frame}
    \frametitle{Linearization}
    \framesubtitle{to preserve Gaussianity}

\begin{columns}
\column{.55\textwidth}%
    \begin{equation}
        \mathrm T_{s,r}[v] \approx \int_C \frac 1{\mu_V(r)} - \frac{V(r)}{\mu_V(r)^2} \mathrm d r
    \end{equation}
    \begin{description}[leftmargin=!, labelwidth=1cm]
        \item [Taylor expansion] 1\textsuperscript{st} order
        \item [point of expansion] $\mu_V$
    \end{description}
    \medskip

    \begin{block}{Approximated correlations}
    \begin{equation}
        \Cov{V}{\mathrm T_{s,r}[V]\,} \approx -\int_C \frac {K_V(\cdot,r)}{\mu_V(r)^2} \mathrm d r
    \end{equation}
    \end{block}

    \begin{block}{Approximated covariance}
    \setlength\abovedisplayskip{0pt}
    \begin{equation}
        \Cov{\mathrm T_{s,r}[V]}{\mathrm T_{\acute s, \acute r}[V]\,} \approx  \int_C \int_{\acute C} \frac{K_V(r,\acute r)}{\mu_V(r)^2\mu_V(\acute r)^2} \mathrm d r \mathrm d \acute r
    \end{equation}
    \end{block}

\column[T]{.44\textwidth}
    \vspace{-10mm}
    \input{fig_correlation.pgf}

\end{columns}

\end{frame}


\begin{frame}
    \frametitle{Bayesian Posterior Distribution}
    \framesubtitle{Gaussian Process Regression}

\begin{columns}
\column{.55\textwidth}%
    \begin{equation}
        D = T + E \sim \Gauss{\mu_D}{\Sigma_{DD}}
    \end{equation}
    \begin{description}[leftmargin=!, labelwidth=1cm]
        \item [Error model]        $E\sim \Gauss{0}{\sigma_\varepsilon^2 = \SFWepsilon^2 }$
        \item [Prior travel times] $\mu_D = \Mean{T_{s,r}[\mu_V]\,}$
        \item [Covariance matrix]  $\Sigma_{DD} = \Cov TT  + \Identity \sigma_\varepsilon^2$
    \end{description}

    \begin{block}{Conditional mean and covariance}
    \setlength\abovedisplayskip{0pt}
    \begin{alignat}{3}
        \Mean{V|d} &= \mu_V &&+ \Cov VD \Var{D}^{-1} \big( d - \Mean{D} \big)
        \\
        \Var{V|d}  &= K_V   &&- \Cov VD \Var{D}^{-1} \Cov DV
    \end{alignat}
    \end{block}

    \begin{alertblock}{Accommodate non-linearity}
        \begin{itemize}
            \item Single evidence at a time
            \item Correlations and Variances from predecessor
        \end{itemize}
    \end{alertblock}


\column[T]{.44\textwidth}
    \vspace{-10mm}
    \input{fig_discretization.pgf}
    \small
    \begin{align}
        \delta\sphericalangle &\approx \SFWdeltaangle\,^\circ &
        & \leadsto &
        N &= \SFWnpts
    \end{align}

\end{columns}

\end{frame}


\begin{frame}
    \frametitle{Successive Approach}
    \begin{center}
    \movie[height=0.85\textheight, width=1.51\textheight, autostart]{\includegraphics[height=0.85\textheight, width=1.51\textheight]{animation_pst}}{animation.mp4}
    \end{center}
\end{frame}


\begin{frame}
    \frametitle{Conclusions}

    Preliminary results
    \begin{itemize}
        \item Accounts for irregular data distributions
        \item Linearization performs well
        \item Realistic assessment of spatial uncertainties
        \item Computationally simple
        \item All calculus in the function space view \\
              No basis functions, no truncations at a certain degree
        \item Memory intens
    \end{itemize}
    \medskip

    Outlook
    \begin{itemize}
        \item Compare with models providing uncertainties
        \item Use of better suited correlation structures
        \item Succesion order
        \item Dispersion;
            kernel accounting for individual frequency bands
        \item Real data;
            estimation of hyper parameters
    \end{itemize}

\end{frame}

\begin{frame}[c, noframenumbering, plain]
    \hspace{15mm} Follow the project on GitHub \\[1cm]
    \hfill \fbox{\qrcode[height=3cm]{https://github.com/mauimuc/gptt}} \hspace{20mm}
\end{frame}

\end{document}
